(1) 有 \(\mathbf R_s=(x_s,y_s,1)\)，\(\mathbf R_{ss}=(x_{ss},y_{ss},0)\)。小斜率 \(|x_s|,|y_s|\ll1\) 时，
\(|\mathbf R_s|^{-3}=1+O(|\psi_s|^2)\)，叉积横向分量主阶是 \((-y_{ss},x_{ss})\)，于是
\[
x_t=-\beta y_{ss},\quad y_t=\beta x_{ss},\quad
\boxed{\psi_t=i\beta\psi_{ss}}.
\]
这与自由Schrödinger方程数学同构，但 \(\psi\) 是涡丝位移，并非概率振幅。

(2) 代入螺旋波，得到 \(\omega=\beta k^2\)。
相速度 \(v_{\mathrm{ph}}=\omega/k=\beta k\)（\(k\ne0\)），群速度
\(v_g=d\omega/dk=2\beta k\)。正负波数沿相反轴向传播，螺旋手性也翻转，尽管频率非负。不同 \(k\) 的群速度不同，故有限谱宽的包必然展宽；当波长接近涡核半径时，固定 \(\beta\) 的局域诱导模型不再可靠，不能把 \(k\to\infty\) 当成真实涡丝的准确预言。

(3) 原方程及共轭方程相减得
\[
\partial_t|\psi|^2+\partial_s J_0=0,\quad
J_0=2\beta\operatorname{Im}(\psi^*\psi_s).
\]
将相同计算应用于 \(\psi_s\) 得
\[
\partial_t|\psi_s|^2+\partial_s J_1=0,\quad
J_1=2\beta\operatorname{Im}(\psi_s^*\psi_{ss}).
\]
边界总流量为零时 \(N=\int|\psi|^2ds\) 及
\(G=\int|\psi_s|^2ds\) 守恒。涡线弧长小振幅增量
\(\Delta\ell=\frac12G+O(|\psi_s|^4)\)，但流体真实能量还可能包含非局域项，故 \(G\) 并非无条件的三维总动能。

(4) 以Fourier变换求解自由二次色散，令 \(d=1+2i\beta t/\sigma^2\)，得到
\[
\psi(s,t)=\frac{A}{\sqrt d}
e^{ik_0s-i\beta k_0^2t}
\exp\left[-\frac{(s-2\beta k_0t)^2}{2\sigma^2d}\right].
\]
强度中心为 \(s_c=2\beta k_0t\)，方差
\[
\operatorname{Var}(s)=\frac{\sigma^2}{2}
\left(1+\frac{4\beta^2t^2}{\sigma^4}\right).
\]
直接高斯积分给 \(\int|\psi|^2ds=|A|^2\sqrt\pi\,\sigma\)，与第三问独立符合。短时 \(|\beta t|\ll\sigma^2\) 恢复原宽度与主阶群速平移；长时宽度线性增长是色散，不是耗散。
